\subsection{Statistical Experiments and Likelihoods}

A statistical problem starts with observable data and a family of probability laws that could have generated it.  The model is not the data itself; it is the mathematical language in which uncertainty about the data-generating mechanism is expressed.

\begin{definition} \label{def:statistical-experiment-and-model} \hlt{Statistical Experiment and Model}
\begin{enumerate}[label=(\roman*)]
\setlength{\itemsep}{0pt}
\item A \hlt{statistical experiment} is $(\mathcal{X},\mathcal{A},\mathcal{P})$, where $(\mathcal{X},\mathcal{A})$ is measurable and\\
$\mathcal{P}=\{P_\theta:\theta\in\Theta\}$ is a family of probability measures on $(\mathcal{X},\mathcal{A})$.
\item $\Theta$ is the \hlt{parameter space}, $X\in\mathcal{X}$ is the \hlt{sample}, and $\mathcal{P}$ is the \hlt{statistical model}.
\item The model is \hlt{well specified} if the true law $P_0\in\mathcal{P}$.
\item The model is \hlt{misspecified} if $P_0\notin\mathcal{P}$.
\end{enumerate}
\end{definition}

\begin{definition} \label{def:dominated-model-likelihood-score-information} \hlt{Dominated Models, Likelihood, Score, and Information}
\begin{enumerate}[label=(\roman*)]
\setlength{\itemsep}{0pt}
\item A model $\mathcal{P}=\{P_\theta:\theta\in\Theta\}$ is \hlt{dominated} by $\mu$ if $P_\theta\ll\mu$ for all $\theta\in\Theta$; write $p_\theta\defeq \frac{dP_\theta}{d\mu}$.
\item Given $x$, the \hlt{likelihood} and \hlt{log-likelihood} are $L_x(\theta)\defeq p_\theta(x)$ and $\ell_x(\theta)\defeq\log p_\theta(x)$ when $p_\theta(x)>0$.
\item If $\Theta\subseteq\mathbb{R}^d$ and differentiation is valid, the \hlt{score} is $S_\theta(X)\defeq\nabla_\theta\ell_X(\theta)$.
\item The \hlt{Fisher information} is $I(\theta)\defeq\mathbb{E}_\theta[S_\theta(X)S_\theta(X)^\top]$.
\end{enumerate}
\end{definition}

\begin{definition} \label{def:maximum-likelihood-estimator} \hlt{Maximum Likelihood Estimator}\\
For likelihood $L_X(\theta)$ and log-likelihood $\ell_X(\theta)$, a \hlt{maximum likelihood estimator} is any measurable\\
$\widehat{\theta}_n(X)\in\arg\max\limits_{\theta\in\Theta}L_X(\theta)=\arg\max\limits_{\theta\in\Theta}\ell_X(\theta)$; equivalently, $\widehat{\theta}_n\in\arg\min\limits_{\theta\in\Theta}[-\ell_X(\theta)]$.
\end{definition}

\begin{example} \label{ex:bernoulli-and-normal-likelihoods} \hlt{Bernoulli and Normal Likelihoods}\\
Let $X_1,\ldots,X_n$ be independent observations.
\begin{enumerate}[label=(\roman*)]
\setlength{\itemsep}{0pt}
\item If $X_i\sim\operatorname{Bernoulli}(p)$, then $L_x(p)=p^{s}(1-p)^{n-s}$, where $s\defeq\sum\limits_{i=1}^{n}x_i$.
\item If $X_i\sim N(\mu,\sigma^2)$ and $\sigma^2$ is known, then $\ell_x(\mu)=-\frac{n}{2}\log(2\pi\sigma^2)-\frac{1}{2\sigma^2}\sum\limits_{i=1}^{n}(x_i-\mu)^2$, maximized at $\widehat{\mu}=\overline{x}$.
\end{enumerate}
\end{example}

\begin{example} \label{ex:multinomial-cell-probability-likelihood} \hlt{Multinomial Cell Probabilities}\\
Let $X=(X_1,\ldots,X_m)$ have total count $n$.
\begin{enumerate}[label=(\roman*)]
\setlength{\itemsep}{0pt}
\item If $X\sim\operatorname{Multinomial}(n;p_1,\ldots,p_m)$ and $\sum\limits_{j=1}^{m}p_j=1$, then $L_x(p)\propto\prod\limits_{j=1}^{m}p_j^{x_j}$ and $\widehat{p}_j=\frac{x_j}{n}$.
\item If $p_j=p_j(\theta)$, then $L_x(\theta)\propto\prod\limits_{j=1}^{m}p_j(\theta)^{x_j}$ and $\widehat{\theta}\in\arg\max\limits_\theta\sum\limits_{j=1}^{m}x_j\log p_j(\theta)$.
\end{enumerate}
\end{example}

\begin{definition} \label{def:identifiability-and-parametric-functionals} \hlt{Identifiability and Functionals}
\begin{enumerate}[label=(\roman*)]
\setlength{\itemsep}{0pt}
\item The parameterization is \hlt{identifiable} if $P_\theta=P_{\theta'}$ implies $\theta=\theta'$.
\item A \hlt{statistical functional} is a map $T:\mathcal{P}\to\mathcal{T}$, such as $T(P)=\int x\,dP(x)$ or $T(P)=\operatorname{median}(P)$.
\item In a parametric model, a \hlt{target parameter} may be a function $\psi(\theta)$ rather than the full parameter $\theta$.
\end{enumerate}
\hlt{Key Insight}: Estimation should target a law or functional that the experiment can actually distinguish.
\end{definition}

\begin{remark} \label{rem:likelihood-is-relative-evidence} \hlt{Likelihood as Relative Evidence}\\
For fixed data $x$, the likelihood is a function of $\theta$, not a probability distribution on $\Theta$.\\
Ratios $\frac{L_x(\theta_1)}{L_x(\theta_0)}$ compare how well two model laws explain the same observation.
\end{remark}
